\section{Introduction}

Transformers with parameter sharing, often called \emph{looped} or \emph{recurrent} Transformers, have emerged as an efficient and capable alternative to deep non–shared stacks across vision and natural language \cite{dehghani2018universal, lan2019albert, jaegle2021perceiver, dutta2021redesigning, geiping2025scaling}. In particular, looped Transformers in the language modeling setting have shown promising performance on a variety of algorithmic and reasoning tasks \cite{geiping2025scaling, saunshi2024inductive, saunshi2025reasoning,gatmiry2024can}. These models appear to possess an inductive bias toward reasoning: a property sometimes referred to as \emph{latent reasoning}, where they internalize reasoning skills akin to explicit chain-of-thought prompting in large language models. Moreover, studies indicate that such abilities scale gracefully with effective computational depth and number of loops, yielding improved results on reasoning benchmarks \cite{xu2024expressive, saunshi2025reasoning, geiping2025scaling}. However, existing approaches almost always train and evaluate with a fixed number of unrolls. This raises a fundamental question: do looped Transformers truly exploit their computational depth flexibly, and can they be trained to operate effectively under variable compute budgets?



Despite their promise, current looped models remain tied to a single trajectory length. Once trained, their representations collapse when evaluated at shorter or longer depths, since those settings are off-distribution \cite{bae2024relaxed,fan2024looped}. In practice, this means looped models spend the same budget as non-looped iso-FLOP baselines, forfeiting one of their key motivations: flexible compute. We instead consider \emph{budget-conditioned} language modeling: at inference, a user specifies a compute budget $M$, and the model should produce high-quality representations without retraining. While early exiting, routing, and layer dropping have made non-looped Transformers more dynamic \cite{schuster2022confident, fan2019reducing, elhoushi2024layerskip, shazeer2017outrageously, raposo2024mixture}, little has been explored for looped models. Naively transplanting these techniques into \emph{looped} architectures is fragile \cite{bae2025mixture, geiping2025scaling}: repeated passes of the shared block often converge to similar, \emph{stagnant} states. We aim for \emph{elastic depth}: a single looped model that performs well across user-chosen budgets without retraining or late-step degeneracy. The core challenge is to train looped models whose internal trajectories remain \emph{stable across depths}, so that shorter routes do not degrade and longer routes continue to refine rather than collapse.


We present \textbf{LoopFormer}, a shortcut-modulated looped language model that supports elastic depth, maintaining strong performance across a range of inference budgets. Inspired by diffusion models \cite{frans2024one, lu2024simplifying} and neural ODEs \cite{chen2018neural}, we cast iterative representation refinement as a trajectory in representation space: token states evolve from an initial $h_0$ toward a target $h_1$ over a normalized unit-time horizon. Our key insight is to explicitly condition each loop step on the current time $t$ and a step size $\Delta t$ (a “jump”), allowing coarser trajectories to approximate fine-grained ones with fewer steps. Training employs a \emph{shortcut-consistency} objective that aligns the final representations of shorter routes with those of the full route via a stop-gradient target, effectively performing self-distillation within the loop. At inference, this conditioning yields \emph{elastic depth} without retraining: the user selects a budget $M \le L$ (maximum loops) and a step schedule, and performance scales smoothly with compute, rather than collapsing at shorter depths.

Departing from dynamic compute through routing or token halting \cite{dehghani2018universal, bae2025mixture}, we instead introduce the notion of \emph{loop trajectories}, which enable effective and efficient language modeling, thereby giving the model the ability to internalize and refine reasoning and performance as compute increases. Empirically, LoopFormer not only preserves the language modeling abilities reported in prior work, but also outperforms baseline looped models on language tasks and closes much of the gap with iso-FLOP non-looped models. This creates an opportunity to study loop trajectories as a new lens on reasoning in looped models. Moreover, unlike early-exit mechanisms that often collapse to degenerate states, LoopFormer maintains stable refinement: both perplexity and reasoning scale gracefully with the compute budget.


The key contributions of this paper are summarized as follows:
\begin{enumerate}[leftmargin=*,labelindent=0pt,label=\textbf{--}]
    \item We formulate a class of shortcut-modulated looped Transformers for language modeling that conditions each pass on internal time and step sizes. 
    
    \item We introduce a shortcut-consistency training protocol over families of step schedules that enables compute-budgeted inference (elastic depth) without retraining.

    \item We analyze representation dynamics across loop iterations using multiple geometric and information-theoretic metrics (anisotropy, curvature, entropy, CKA), finding that adaptive early-exit looped models exhibit representational collapse, while our shortcut-modulated models maintain evolving, non-degenerate states.
    
    \item We demonstrate consistent performance-per-compute gains in both perplexity and zero-shot reasoning across diverse language benchmarks; ablations isolate the roles of time/step conditioning and offer practical guidance for choosing trajectories under fixed budgets.

\end{enumerate}


